% Part: first-order-logic
% Chapter: completeness
% Section: complete-sets

% Definition of complete consistent sets. Properties of complete sets required
% for completeness proved are in provability.tex in the
% chapter on the proof system used.

\documentclass[../../../include/open-logic-section]{subfiles}

\begin{document}

\iftag{FOL}
      {\olfileid{fol}{com}{ccs}}
      {\olfileid{pl}{com}{ccs}}

\olsection{مجموعات ال\usetoken{P}{sentence} المتسقة الكاملة}

\begin{defn}[المجموعة الكاملة]
\ollabel{def:complete-set} تكون مجموعة~$\Gamma$ من ال!!{sentence}s
\emph{!!{complete}} إذا وفقط إذا تحقق لكل !!{sentence}~$!A$ أن
$!A \in \Gamma$ أو $\lnot !A \in \Gamma$.
\end{defn}

\begin{explain}
مجموعات الجمل ال!!^{complete} تفصل في كل سؤال: فأيّ !!{sentence}~$!A$
أخذنا، وجدنا في $\Gamma$ حكمًا بصدق $!A$ أو كذبها. وليست فائدة
المجموعات ال!!{complete} مقصورة على برهان الاكتمال؛ فإن النظرية
ال!!{complete} إذا كانت قابلة للأكسمة كانت قابلة للقرار.
\end{explain}

\begin{explain}
تقوم المجموعات المتسقة ال!!^{complete} بدور أساس في البرهان؛ إذ سنضمن
أن كل مجموعة متسقة من ال!!{sentence}s~$\Gamma$ تقع ضمن مجموعة متسقة
!!a{complete}~$\Gamma^*$. ففي المجموعة المتسقة !!^a{complete}،
لكل !!{sentence}~$!A$، إما $!A$ وإما نفيها $\lnot !A$، ولا تجتمعان.
ويشمل هذا \iftag{FOL}{ال!!{sentence}s الذرية}{المتغيرات القضوية}؛
ومن ثم تتيح لنا المجموعة المتسقة !!a{complete}\iftag{FOL}{، في لغة
زِيدت فيها !!{constant}s على وجه مناسب،}{} بناء
\iftag{FOL}{!!a{structure}}{!!a{valuation}} نعيّن فيها
\iftag{FOL}{تفسير ال!!{predicate}s}{قيم صدق المتغيرات القضوية}
بحسب ما في~$\Gamma^*$ من
\iftag{FOL}{ال!!{sentence}s الذرية}{ال!!{propositional variable}s}.
ثم نثبت أن هذه ال\iftag{FOL}{!!{structure}}{!!{valuation}} ترضي
ال!!{sentence}s كلها في~$\Gamma^*$، وبذلك ترضي ما في~$\Gamma$.
ولا يتم هذا الإثبات إلا بخواص منها أن $\lnot !A \in \Gamma^*$ إذا
وفقط إذا كان $!A \notin \Gamma^*$، وأن $(!A \lor !B) \in \Gamma^*$
إذا وفقط إذا كان $!A \in \Gamma^*$ أو $!B \in \Gamma^*$، ونظائر ذلك.
\end{explain}

وسنجري في البراهين الآتية على انعكاسية العلاقة $\Proves$ ورتابتها
وتعديها، وكثيرًا ما نكتفي باستعمالها من غير تنبيه (انظر
\iftag{FOL}{%
  \tagrefs{prfSC/{fol:seq:ptn:sec},prfND/{fol:ntd:ptn:sec},prfAX/{fol:axd:ptn:sec},prfTab/{fol:tab:ptn:sec}}}{%
  \tagrefs{prfSC/{pl:seq:ptn:sec},prfND/{pl:ntd:ptn:sec},prfAX/{pl:axd:ptn:sec},prfTab/{pl:tab:ptn:sec}}}).

\begin{prop}
\ollabel{prop:ccs}
افترض أن $\Gamma$ !!{complete} ومتسقة. عندئذ:
\begin{enumerate}
\item \ollabel{prop:ccs-prov-in} إذا كان $\Gamma \Proves !A$، فإن
  $!A \in \Gamma$.

\tagitem{prvAnd}{\ollabel{prop:ccs-and} $!A \land !B \in \Gamma$
  إذا وفقط إذا كان كل من $!A \in \Gamma$ و$!B \in \Gamma$.}{}

\tagitem{prvOr}{\ollabel{prop:ccs-or} $!A \lor !B \in \Gamma$ إذا
  وفقط إذا كان $!A \in \Gamma$ أو $!B \in \Gamma$.}{}

\tagitem{prvIf}{\ollabel{prop:ccs-if} $!A \lif !B \in \Gamma$ إذا
  وفقط إذا كان $!A \notin \Gamma$ أو $!B \in \Gamma$.}{}
\end{enumerate}
\end{prop}

\begin{proof}
فرضنا في جميع أجزاء البرهان أن $\Gamma$ !!{complete} مع اتساقها.
\begin{enumerate}
\item إذا كان $\Gamma \Proves !A$، فإن $!A \in \Gamma$.

ليكن $\Gamma \Proves !A$. فإن فرضنا $!A \notin
\Gamma$، اقتضى كون $\Gamma$ !!{complete} أن $\lnot !A \in \Gamma$.
فتفيد
\iftag{FOL}{%
  \tagrefs{prfSC/{fol:seq:prv:prop:explicit-inc},
    prfND/{fol:ntd:prv:prop:explicit-inc},
    prfAX/{fol:axd:prv:prop:explicit-inc},
    prfTab/{fol:tab:prv:prop:explicit-inc}}}{%
  \tagrefs{prfSC/{pl:seq:prv:prop:explicit-inc},
    prfND/{pl:ntd:prv:prop:explicit-inc},
    prfAX/{pl:axd:prv:prop:explicit-inc},
    prfTab/{pl:tab:prv:prop:explicit-inc}}},
أن $\Gamma$ غير متسقة، خلاف اتساق $\Gamma$ المفروض. فبطل
$!A \notin \Gamma$، وثبت $!A \in \Gamma$.

\tagitem{defAnd}{}{%
\iftag{probAnd}{تمرين.}{%
$!A \land !B \in \Gamma$ إذا وفقط إذا كان كل من $!A \in \Gamma$ و
$!B \in \Gamma$:

أما اللزوم من اليسار إلى اليمين، فليكن $!A \land !B \in \Gamma$.
بمقتضى
\iftag{FOL}{%
  \tagrefs{prfSC/{fol:seq:ppr:prop:provability-land},
    prfND/{fol:ntd:ppr:prop:provability-land},
    prfAX/{fol:axd:ppr:prop:provability-land},
    prfTab/{fol:tab:ppr:prop:provability-land}}}{%
  \tagrefs{prfSC/{pl:seq:ppr:prop:provability-land},
    prfND/{pl:ntd:ppr:prop:provability-land},
    prfAX/{pl:axd:ppr:prop:provability-land},
    prfTab/{pl:tab:ppr:prop:provability-land}}}، البند~(1)،
يلزم $\Gamma \Proves !A$ و$\Gamma \Proves !B$؛ فتنقلنا
\olref{prop:ccs-prov-in} إلى $!A \in \Gamma$ و$!B \in \Gamma$،
وهو المطلوب.

وأما العكس، فإذا كان $!A \in \Gamma$ و$!B \in \Gamma$، أفاد
\iftag{FOL}{%
  \tagrefs{prfSC/{fol:seq:ppr:prop:provability-land},
    prfND/{fol:ntd:ppr:prop:provability-land},
    prfAX/{fol:axd:ppr:prop:provability-land},
    prfTab/{fol:tab:ppr:prop:provability-land}}}{%
  \tagrefs{prfSC/{pl:seq:ppr:prop:provability-land},
    prfND/{pl:ntd:ppr:prop:provability-land},
    prfAX/{pl:axd:ppr:prop:provability-land},
    prfTab/{pl:tab:ppr:prop:provability-land}}}، البند~(2)،
أن $\Gamma \Proves !A \land !B$؛ وبـ\olref{prop:ccs-prov-in}
يكون $!A \land !B \in \Gamma$.}}

\tagitem{defOr}{}{%
\iftag{probOr}{تمرين.}{%
نبدأ بإثبات أن $!A \lor !B \in \Gamma$ يقتضي $!A \in \Gamma$
أو $!B \in \Gamma$. فلو كان $!A \lor !B \in \Gamma$ مع
$!A \notin \Gamma$ و$!B \notin \Gamma$، للزم، لأن $\Gamma$
!!{complete}، أن $\lnot !A \in \Gamma$ و$\lnot !B \in \Gamma$.
فتفيد
\iftag{FOL}{%
  \tagrefs{prfSC/{fol:seq:ppr:prop:provability-lor},
    prfND/{fol:ntd:ppr:prop:provability-lor},
    prfAX/{fol:axd:ppr:prop:provability-lor},
    prfTab/{fol:tab:ppr:prop:provability-lor}}}{%
  \tagrefs{prfSC/{pl:seq:ppr:prop:provability-lor},
    prfND/{pl:ntd:ppr:prop:provability-lor},
    prfAX/{pl:axd:ppr:prop:provability-lor},
    prfTab/{pl:tab:ppr:prop:provability-lor}}}،
البند (1)، عدم اتساق $\Gamma$، وهو محال. فلا بد من
$!A \in \Gamma$ أو $!B \in \Gamma$.

وأما العكس، فليكن $!A \in \Gamma$ أو $!B \in \Gamma$؛ وعندئذ يفيد
\iftag{FOL}{%
  \tagrefs{prfSC/{fol:seq:ppr:prop:provability-lor},
    prfND/{fol:ntd:ppr:prop:provability-lor},
    prfAX/{fol:axd:ppr:prop:provability-lor},
    prfTab/{fol:tab:ppr:prop:provability-lor}}}{%
  \tagrefs{prfSC/{pl:seq:ppr:prop:provability-lor},
    prfND/{pl:ntd:ppr:prop:provability-lor},
    prfAX/{pl:axd:ppr:prop:provability-lor},
    prfTab/{pl:tab:ppr:prop:provability-lor}}}، البند (2)،
أن $\Gamma \Proves !A \lor !B$، فبمقتضى \olref{prop:ccs-prov-in}
ثبت $!A \lor !B \in \Gamma$، وتم هذا الاتجاه.}}

\tagitem{defIf}{}{%
\iftag{probIf}{تمرين.}{%
لإثبات الاتجاه الأول، ليكن $!A \lif !B \in \Gamma$، ولنفرض نقيض
المطلوب: $!A \in \Gamma$ و$!B \notin \Gamma$. فيجتمع
$!A \lif !B \in \Gamma$ مع $!A \in \Gamma$. وعليه يفيد
\iftag{FOL}{%
  \tagrefs{prfSC/{fol:seq:ppr:prop:provability-lif},
    prfND/{fol:ntd:ppr:prop:provability-lif},
    prfAX/{fol:axd:ppr:prop:provability-lif},
    prfTab/{fol:tab:ppr:prop:provability-lif}}}{%
  \tagrefs{prfSC/{pl:seq:ppr:prop:provability-lif},
    prfND/{pl:ntd:ppr:prop:provability-lif},
    prfAX/{pl:axd:ppr:prop:provability-lif},
    prfTab/{pl:tab:ppr:prop:provability-lif}}}، البند (1)،
أن $\Gamma \Proves !B$؛ ومن \olref{prop:ccs-prov-in} يلزم
$!B \in \Gamma$، مناقضًا فرض $!B \notin \Gamma$.

وللعكس حالتان. أما حالة $!A \notin \Gamma$، فكون
$\Gamma$ !!{complete} يقتضي $\lnot !A \in \Gamma$؛ فيفيد
\iftag{FOL}{%
  \tagrefs{prfSC/{fol:seq:ppr:prop:provability-lif},
    prfND/{fol:ntd:ppr:prop:provability-lif},
    prfAX/{fol:axd:ppr:prop:provability-lif},
    prfTab/{fol:tab:ppr:prop:provability-lif}}}{%
  \tagrefs{prfSC/{pl:seq:ppr:prop:provability-lif},
    prfND/{pl:ntd:ppr:prop:provability-lif},
    prfAX/{pl:axd:ppr:prop:provability-lif},
    prfTab/{pl:tab:ppr:prop:provability-lif}}}، البند (2)،
أن $\Gamma \Proves !A \lif !B$؛ وبالرجوع إلى
\olref{prop:ccs-prov-in} يثبت $!A \lif !B \in \Gamma$.

وأما حالة $!B \in \Gamma$، فنستعمل
\iftag{FOL}{%
  \tagrefs{prfSC/{fol:seq:ppr:prop:provability-lif},
    prfND/{fol:ntd:ppr:prop:provability-lif},
    prfTab/{fol:tab:ppr:prop:provability-lif},
    prfAX/{fol:axd:ppr:prop:provability-lif}}}{%
  \tagrefs{prfSC/{pl:seq:ppr:prop:provability-lif},
    prfND/{pl:ntd:ppr:prop:provability-lif},
    prfTab/{pl:tab:ppr:prop:provability-lif},
    prfAX/{pl:axd:ppr:prop:provability-lif}}}،
البند (2) كذلك، فنجد $\Gamma \Proves !A \lif !B$؛ وتثبت
\olref{prop:ccs-prov-in} أن $!A \lif !B \in \Gamma$.}}
\end{enumerate}
\end{proof}

\tagprob{FOL}
\begin{prob}
أكمل برهان \olref[fol][com][ccs]{prop:ccs}.
\end{prob}
\tagendprob

\tagprob{notFOL}
\begin{prob}
أكمل برهان \olref[pl][com][ccs]{prop:ccs}.
\end{prob}
\tagendprob

\end{document}
